\documentclass[11pt,a4paper]{report}

%% ----------------------------- packages
\usepackage[margin=1in]{geometry}
\usepackage{amsmath, amssymb, amsthm, mathtools}
\usepackage{bm}
\usepackage{graphicx}
\usepackage{hyperref}
\usepackage{booktabs}
\usepackage{xcolor}
\usepackage{listings}

\hypersetup{
  colorlinks=true,
  linkcolor=blue!50!black,
  urlcolor=blue!60!black,
  citecolor=teal!60!black,
}

\lstset{
  basicstyle=\ttfamily\small,
  frame=single,
  language=Python,
  keywordstyle=\color{blue!60!black},
  commentstyle=\color{green!40!black}\itshape,
  showstringspaces=false,
}

\newcommand{\bs}[1]{\boldsymbol{#1}}
\newcommand{\Tr}{\operatorname{Tr}}
\newcommand{\dev}{\operatorname{dev}}

\title{femsolver --- Theory Manual}
\author{femsolver contributors}
\date{\today}

\begin{document}
\maketitle
\tableofcontents
\clearpage

%% ============================================================
\chapter{Introduction}
\label{ch:introduction}

\texttt{femsolver} implements the standard structural-engineering
finite-element formulation: displacement-based interpolation,
isoparametric mapping, Gauss quadrature, and assembly into a global
system $\bs K \bs u = \bs f$ that is solved by a direct sparse
factorisation (or iteratively).  The library is organised in layers
(see the Sphinx site for a graphical view): nodes, materials,
sections, elements, constraints, assembly, and analysis.

This manual is the \emph{theory companion} to the docstrings, which
contain code-level detail.  Numbered references below match the
\texttt{src/femsolver/} module layout.

%% ============================================================
\chapter{Linear analysis}
\label{ch:linear}

\section{Equilibrium}

A linear elastic structure satisfies
\begin{equation}
  \bs K \bs u \;=\; \bs f,
  \label{eq:linear-static}
\end{equation}
with $\bs K = \mathop{\mathrm{A}}_e \bs k_e$ the global stiffness
matrix assembled element-wise, $\bs u$ the nodal-displacement
vector, and $\bs f$ the external nodal load.

\section{Element stiffness}

For an element with shape functions $\bs N(\bs \xi)$ over its
natural domain and constitutive matrix $\bs D$,
\begin{equation}
  \bs k_e \;=\; \int_{\Omega_e} \bs B^{\mathsf T} \bs D \bs B \, \mathrm{d}\Omega,
  \qquad
  \bs B \;=\; \frac{\partial \bs N}{\partial \bs x}.
  \label{eq:k-element}
\end{equation}

\section{Mass matrix}

Consistent mass: $\bs m_e = \int_{\Omega_e} \rho \bs N^{\mathsf T}
\bs N \, \mathrm{d}\Omega$.  Lumped: row-sum of the consistent
matrix.

\section{Assembly}

Element-local DOFs are scattered into the global DOF numbering via
the connectivity matrix.  Multi-point constraints (rigid diaphragms,
equal-DOF, rigid links) are imposed by the transformation-handler
method: a sparse $\bs T$ matrix maps the constrained system to the
master DOFs, and the global system becomes $(\bs T^{\mathsf T} \bs K
\bs T) \bs u_m = \bs T^{\mathsf T} \bs f$.

%% ============================================================
\chapter{Nonlinear static analysis}
\label{ch:nonlinear}

\section{Equilibrium iteration}

For a nonlinear structure
\begin{equation}
  \bs R(\bs u) \;=\; \bs f_{\mathrm{int}}(\bs u) - \bs f_{\mathrm{ext}}(\lambda) \;=\; \bs 0,
\end{equation}
solved by Newton-Raphson:
\begin{equation}
  \bs K_T^{(k)} \Delta \bs u^{(k+1)} \;=\; - \bs R^{(k)}, \qquad
  \bs u^{(k+1)} = \bs u^{(k)} + \Delta \bs u^{(k+1)}.
\end{equation}

\section{Integrators}

\paragraph{Load control.}  $\lambda_{n+1} = \lambda_n + \Delta\lambda$.

\paragraph{Displacement control.}  A target DOF is incremented by a
fixed $\mathrm{d} u$ and $\lambda$ becomes an unknown; see Crisfield
1991.

\paragraph{Arc-length.}  Crisfield's spherical arc-length: an extra
constraint $\Delta \bs u^{\mathsf T} \Delta \bs u + \psi^2
(\Delta\lambda)^2 (\bs f^{\mathsf T} \bs f) = \ell^2$ permits
snap-through past limit points.

\section{Convergence tests}

\begin{itemize}
  \item \texttt{NormDispIncr}: $\| \Delta \bs u \| / \| \bs u \| < \mathrm{tol}$.
  \item \texttt{NormUnbalance}: $\| \bs R \| < \mathrm{tol}$.
  \item \texttt{EnergyIncr}: $\frac{1}{2} \Delta \bs u^{\mathsf T} \bs R < \mathrm{tol}$.
\end{itemize}

%% ============================================================
\chapter{Constitutive models}
\label{ch:constitutive}

\section{$J_2$ plasticity}

Yield surface $f = \sqrt{3 J_2} - \sigma_y(\alpha)$ with $J_2 =
\tfrac{1}{2} \dev(\bs\sigma) : \dev(\bs\sigma)$.  Return-mapping
follows Simo \& Hughes 1998 (radial return).  Hardening can be
isotropic ($\sigma_y = \sigma_{y0} + H \alpha$), kinematic (back-stress
$\bs q$), or mixed.

\section{Drucker-Prager}

$f = \alpha I_1 + \sqrt{J_2} - k$ with $I_1 = \Tr(\bs\sigma)$ and
material constants $\alpha, k$ tied to a Mohr-Coulomb fit (friction
angle $\phi$, cohesion $c$).  Non-associated flow with dilation angle
$\psi$.

\section{Uniaxial library}

Each uniaxial material implements
\begin{equation}
  \sigma_{n+1}, \; E_{t,n+1} \;=\; \mathtt{get\_response}(\varepsilon_{n+1}),
\end{equation}
plus \texttt{commit\_state} / \texttt{revert\_state} for path-dependent
models.  The catalogue includes:

\begin{itemize}
  \item \textbf{Bilinear} kinematic / isotropic, with back-stress.
  \item \textbf{Menegotto-Pinto} (evolving R) for steel rebar.
  \item \textbf{Kent-Park}, \textbf{Mander} concrete with confinement
        and Chang-Mander cyclic damage.
  \item \textbf{Hysteretic} with pinching + degradation.
  \item \textbf{IMK} (Ibarra-Medina-Krawinkler) with post-cap softening
        for collapse-capacity work.
  \item \textbf{BRB} asymmetric hardening for buckling-restrained braces.
  \item \textbf{Takeda}, \textbf{Pivot} for RC plastic hinges.
  \item \textbf{Gap}, \textbf{Elastic} as utility models.
\end{itemize}

%% ============================================================
\chapter{Sections and elements}
\label{ch:elements}

\section{Beam-column kinematics}

\paragraph{Bernoulli (displacement-based).}  Linear axial $+$ Hermite
cubic transverse shape functions; integrate over the element axis
with 2-point Gauss-Legendre quadrature.

\paragraph{Corotational.}  Rigid-body motion is extracted to a local
chord frame; basic deformations are axial extension and end
rotations.  The element-internal force in the global frame is
$\bs f = \bs T^{\mathsf T} \bs f_b$ with $\bs T(\bs u_n)$
nonlinear in nodal positions.

\paragraph{Force-based.}  Neuenhofer-Filippou 1997 state determination
in flexibility form; iteratively enforces section equilibrium at the
integration points.

\section{Plate / shell}

MITC4 (Bathe-Dvorkin assumed strains for transverse shear), MITC9,
DKMQ4, ShellTri3, and the discrete-Kirchhoff DKT3 triangle.  All ship
with full geometric-stiffness terms for buckling.

\section{Fiber sections}

The cross-section is discretised into uniaxial fibers $f_i$ with
position $(y_i, z_i)$, area $A_i$, and material law $\sigma_i =
\mathtt{material}_i(\varepsilon_i)$.  Plane sections give
\begin{equation}
  \varepsilon_i \;=\; \varepsilon_a - y_i \kappa_z + z_i \kappa_y,
\end{equation}
and section forces aggregate as
\begin{equation}
  \bs s = (N, M_z, M_y, T)^{\mathsf T}, \qquad
  N = \sum \sigma_i A_i, \quad M_z = - \sum y_i \sigma_i A_i, \;\text{etc.}
\end{equation}

%% ============================================================
\chapter{Dynamics}
\label{ch:dynamics}

\section{Equation of motion}

\begin{equation}
  \bs M \ddot{\bs u} + \bs C \dot{\bs u} + \bs f_{\mathrm{int}}(\bs u)
  \;=\; \bs f_{\mathrm{ext}}(t).
\end{equation}

Linear Rayleigh damping: $\bs C = \alpha_M \bs M + \alpha_K \bs K$
with $\alpha_M, \alpha_K$ tuned to give a specified $\zeta_i, \zeta_j$
at two target frequencies.

\section{Newmark family}

Average-acceleration ($\beta = 1/4, \gamma = 1/2$) is
unconditionally stable.  HHT-$\alpha$ and Generalised-$\alpha$ add
controlled numerical dissipation at high frequencies (Chung-Hulbert
1993).  Central difference is explicit, conditionally stable.

\section{Modal analysis}

Generalised eigenproblem $\bs K \bs\phi_i = \omega_i^2 \bs M \bs\phi_i$
solved by \texttt{scipy.sparse.linalg.eigsh} (shift-invert).  Mass-
orthonormalised modes feed response-spectrum, MPA, and modal
participation factors.

\section{Response spectrum}

Peak modal response is $u_i = \Gamma_i S_a(T_i, \zeta) /
\omega_i^2 \bs\phi_i$, combined by SRSS or CQC.  Implementation
matches Chopra 2017 Eqs. 13.7.\,4 -- 13.7.\,11.

%% ============================================================
\chapter{Performance-based earthquake engineering}
\label{ch:pbe}

\section{Pipeline}

Hazard $\to$ records $\to$ IDA $\to$ collapse $\to$ fragility $\to$
component damage (FEMA P-58) $\to$ loss curve.  Implementation is
in \texttt{femsolver.performance.ida}, \texttt{ida\_collapse},
\texttt{fragility}, \texttt{record\_scaling}, \texttt{cms},
\texttt{p58}.

\section{IDA}

Per Vamvatsikos-Cornell 2002: each record is scaled to a sequence of
intensity-measure (IM) levels; a full nonlinear time-history analysis
runs at each (record, IM) cell.  Collapse is flagged by the union of
three criteria: hard EDP limit (drift $> 10\%$), NLTHA non-convergence,
or slope-flatlining of the IDA curve.

\section{Fragility}

Lognormal CDF $P(\mathrm{collapse} \mid \mathrm{IM}) = \Phi\!\left(
\frac{\ln(\mathrm{IM} / \theta)}{\beta} \right)$.  Fitting is by
method-of-moments (all records collapsed) or by censored MLE
(some records did not collapse within the swept range; FEMA P-695
Appendix F).

\section{Conditional Mean Spectrum}

Baker 2011: the spectral shape of records conditioned on $\mathrm{Sa}(T^*)
= \mathrm{Sa}_{\mathrm{target}}$ is
\begin{equation}
  \ln \mathrm{CMS}(T_i) \;=\; \mu_{\ln \mathrm{Sa}}(T_i)
  \;+\; \rho(T_i, T^*) \, \epsilon^* \, \sigma_{\ln \mathrm{Sa}}(T_i),
\end{equation}
where $\rho$ is the Baker-Jayaram 2008 empirical inter-period
correlation.  CMS replaces the conservative uniform-hazard spectrum
in modern PBSE record selection.

\section{FEMA P-58 component damage}

Each component has a set of damage states with fragility
$P(\mathrm{DS} \geq \mathrm{DS}_i \mid \mathrm{EDP}) = \Phi(\ln(
\mathrm{EDP} / \theta_i) / \beta_i)$ and repair-cost lognormal
$(\mathrm{cost\_median}, \mathrm{cost\_beta})$.  Expected loss:
\begin{equation}
  E[L] \;=\; Q \sum_i P(\mathrm{DS} = \mathrm{DS}_i \mid \mathrm{EDP}) \,
  E[\mathrm{cost}_i],
\end{equation}
with $E[\mathrm{cost}_i] = c_i \exp(\beta_i^2 / 2)$ (lognormal mean).
Monte-Carlo realisations sample the damage state and consequence per
component group.

%% ============================================================
\chapter{Design and detailing}
\label{ch:design}

\section{Concrete (ACI 318, IS 456)}

Limit-state flexure with rectangular stress block: ACI uses Whitney
$0.85 f_c'$ over depth $a = \beta_1 c$; IS 456 uses $0.36 f_{ck}$ over
$x_u$ with steel at $0.87 f_y$.  Both reduce to a quadratic in the
neutral-axis depth and yield closed-form $A_{st}$ requirements.

\section{Steel (AISC 360, IS 800)}

LRFD member capacities under axial, flexure, shear, and combined
loads.  Perry-Robertson buckling curves (a/b/c/d) and LTB reductions
follow EC3 / IS 800 form; AISC uses $\phi$-factors.

\section{Seismic detailing}

Strong-column / weak-beam (ACI 18.7.3 / AISC 341 E3.4a / IS 13920
Cl. 7.2.1).  Beam capacity shear (ACI 18.6.5 / AISC 341 E3.4b /
IS 13920 Cl. 6.3.3).  Confinement (ACI 18.7.5 / IS 13920 Cl. 7.6).

\section{Connections}

Krawinkler panel-zone bilinear shear backbone with the AISC J10 form;
AISC 358 RBS reduced-section design with the $C_{pr}$ probable
strength factor; Richard-Abbott four-parameter PR connections.

%% ============================================================
\chapter{Bridges}
\label{ch:bridges}

\section{Influence lines}

Closed-form for SS spans; numerical Müller-Breslau for continuous
beams.  Moving-load envelopes for AASHTO HL-93 (truck $+$ tandem $+$
lane) and IRC Class A / 70R.

\section{PT tendons}

Friction loss $P(x) / P_0 = \exp\!\left(-(\mu \alpha + k x)\right)$
along a parabolic-drape profile, plus anchorage-slip redistribution
$l_a = \sqrt{\Delta_a E_s A_{ps} / p_{slope}}$.

\section{Creep, shrinkage, relaxation}

CEB-FIP MC 2010 creep $\phi(t, t_0)$ and shrinkage $\varepsilon_{cs}(t)$,
plus low-relaxation strand loss
$\Delta f_p / f_{pi} = (\log_{10}(t/t_i) / 45)(f_{pi}/f_{py} - 0.55)$
for strand.  Aggregated by \texttt{prestress\_long\_term\_loss}.

%% ============================================================
\chapter{Verification and validation}
\label{ch:vv}

The shipped V\&V suite (Phase~35) checks closed-form references for
beams, plates, shells, solids, eigen, buckling, and plasticity.  See
\texttt{src/femsolver/benchmarks/} and the
\texttt{examples/46\_vnv\_report.py} capstone.

A summary of the current pass status is rendered to CSV via
\texttt{export\_csv(results, \dots)} and committed as a regression
artefact.

%% ============================================================
\chapter*{References}
\addcontentsline{toc}{chapter}{References}

\begin{itemize}
  \item Crisfield, M.A. (1991/1997). \emph{Non-linear Finite Element
        Analysis of Solids and Structures}. Wiley.
  \item Simo, J.C. \& Hughes, T.J.R. (1998). \emph{Computational
        Inelasticity}. Springer.
  \item Chopra, A.K. (2017). \emph{Dynamics of Structures}, 5e.
        Pearson.
  \item Vamvatsikos, D. \& Cornell, C.A. (2002). ``Incremental dynamic
        analysis.'' \emph{Earthquake Eng.\ \& Struct.\ Dyn.}, 31(3).
  \item Baker, J.W. (2011). ``Conditional Mean Spectrum.''
        \emph{J.\ Struct.\ Eng.}, 137(3).
  \item Baker, J.W.\ \& Jayaram, N. (2008). ``Correlation of spectral
        acceleration values from NGA ground-motion models.''
        \emph{Earthquake Spectra}, 24(1).
  \item FEMA P-58 (2018). ``Seismic Performance Assessment of
        Buildings.''
  \item ACI 318-19, AISC 360-22, ASCE 7-22.
  \item IS 456:2000, IS 800:2007, IS 1893 (Part 1):2016, IS 13920:2016.
  \item AASHTO LRFD Bridge Design Specifications, 9e (2020).
  \item CEB-FIP Model Code (2010).
\end{itemize}

\end{document}
